\documentclass[../../../include/open-logic-section]{subfiles}

\begin{document}

\olfileid{his}{set}{infinitesimals}
\olsection{无穷小与微分}

Newton 和 Leibniz 在 17 世纪末各自独立地发现了微积分。微积分的一个特别重要的应用是\emph{微分}。粗略地说，微分旨在给出函数在某一点上的“变化率”或梯度的概念。

我们可以用一种直观的方式来解释这个想法。考虑函数 $f(x) = \nicefrac{x^2}{4} + \nicefrac{1}{2}$，如下图中黑色曲线所示：
\begin{center}
	\begin{tikzpicture}[scale=1]
	\draw[->, gray] (-1,0) -- (4.25,0) node[right] {$x$};
	\draw[->, gray] (0,-0.25) -- (0,5.25) node[above] {$f(x)$};
	
	\foreach \x/\xtext in {1/1, 2/2, 3/3, 4/4}
	\draw[shift={(\x,0)}, gray] (0pt,2pt) -- (0pt,-2pt) node[below] {$\xtext$};
	
	\foreach \y/\ytext in {1/1, 2/2, 3/3, 4/4, 5/5}
	\draw[shift={(0,\y)}, gray] (2pt,0pt) -- (-2pt,0pt) node[left] {$\ytext$};
	
	\draw[oldiagcolorC] (0.5, 9/16) -- (3.5, 57/16) -- (3.5, 9/16)--cycle;
	\draw[oldiagcolorD] (.5, 9/16) -- (2.5, 33/16) -- (2.5, 9/16) -- cycle;
	\draw[oldiagcolorE] (.5, 9/16) -- (1.5, 17/16) -- (1.5, 9/16) -- cycle;
	\draw[thick] (-1,.75) parabola bend (0,0.5) (4,4.5);
	\end{tikzpicture}
\end{center}
假设我们想求函数在 $c = \nicefrac{1}{2}$ 处的梯度。我们可以先画一个斜边近似于该点梯度的三角形，比如上图中的红色三角形。当三角形的底边长度为 $\beta$ 时，它的高是 $f(\nicefrac{1}{2}+\beta) - f(\nicefrac{1}{2})$，因此斜边的梯度为：
\[
\frac{f(\nicefrac{1}{2}+\beta) - f(\nicefrac{1}{2})}{\beta}.
\]
因此，底边长为~$3$ 的红色三角形，其梯度恰好为~$1$。底边长为~$2$ 的蓝色三角形更小，其斜边给出了一个更好的近似；它的梯度是 $\nicefrac{3}{4}$。底边长为~$1$ 的绿色三角形更小，给出的近似更优；其梯度为 $\nicefrac{1}{2}$。

越来越小的三角形给出越来越好的近似。因此我们或许可以这样说：底边长度为\emph{无穷小}的三角形的斜边，给出了 $c = \nicefrac{1}{2}$ 处的梯度本身。这样一来，我们将得到函数 $f$ 在点 $c$ 处的（一阶）导数公式：
\[
{f'}(c) = \frac{f(c+\beta) - f(c)}{\beta} \text{ 其中 $\beta$ 是无穷小。}
\]
而这大致就是 Newton 和 Leibniz 所说的。

然而，既然他们这样说了，我们就必须问他们：\emph{无穷小}究竟是什么？这引发了一个严重的两难困境。如果 $\beta = 0$，那么 $f'$ 就无法定义了，因为它涉及除以 $0$。但如果 $\beta > 0$，那么我们得到的只是对梯度的\emph{近似}，而非梯度本身。

这并非一种时代错置的担忧。以下是 Berkeley 对 Newton 追随者的批评：
\begin{quote}
我承认，符号可以被用来指代任何事物或者空无：因此，在最初的记法 $c + \beta$ 中，$\beta$ 可能意指一个增量，也可能意指空无。但无论你让它表示哪一种，你的论证就必须与该含义保持一致，而不能基于双重含义推进：那样做显然是诡辩。(\citealt[\S{}XIII]{Berkeley1734}，变量已改为与上文匹配)
\end{quote}
为了替无穷小微积分辩护以回应 Berkeley，有人可能会辩称，关于“无穷小”的说法仅仅是一种比喻。或许会说，只要取一个\emph{极小}的三角形，就能得到切线的一个\emph{足够好}的近似。Berkeley 对此也有反驳：虽然这对工程学来说可能足够好，但它损害了数学的\emph{地位}，因为
\begin{quote}
我们被告知：\emph{in rebus mathematicis errores qu\`{a}m minimi non sunt contemnendi}. [在数学中，最微小的误差也不容忽视。] \citep[\S{}IX]{Berkeley1734}
\end{quote}

这段斜体字几乎是逐字引用自 Newton 本人的《曲线求积术》(1704)。

Berkeley 的哲学反驳极为犀利。尽管如此，微积分是一项极为成功的事业，数学家们在继续使用它的过程中并未因此出错。

\end{document}